\histitem{紧算子}

不到三年之后，里斯在其关于泛函方程的论文 [60] 中大大强化了这些评注；该文以历经一个多世纪仍未减损的清晰度，发展了巴拿赫空间上紧算子谱理论的要点。这篇文本宣称的目标是，通过考虑赋有一致收敛拓扑的空间 \(E = \mathcal{C}(I)\) ，对紧区间 I 上的方程 (1) 重新表述弗雷德霍姆理论。但正如我们在早先的一条注记中已经指出的（ÉHM, p.~268），里斯清楚地意识到他的结果适用于任何巴拿赫空间——这一概念要到下一个十年才被引入。

受弗雷歇关于一般拓扑的思想启发，里斯这时考虑 E 的把每个有界序列映为相对紧序列的自同态。与希尔伯特一样，他仍称其为“全连续的”，并指出在 \(\mathrm{E} = \ell_{\mathbf{C}}^{2}(\mathbf{N})\) 的情形，这一新概念等价于他在 1913 年的书中 \(^{(4)}\) 用过的全连续性概念。

里斯研究形如 \(B = 1_{E} - A\) 的自同态，其中 A 是 E 的全连续自同态，并从“局部紧的赋范向量空间必是有限维的”这一——似乎正是为该场合而发现的——事实导出它们的所有谱性质。他容易地证明 B 的核是有限维的，且它的像是闭的、余维数有限。如同 E. 韦尔 [87] 在有限维中所做的那样，他接着考虑迭代 \(B^{k}\) 。他证明它们的核所成的序列与像所成的序列都是平稳的，然后引入我们今天所谓的 B 的零空间与余零空间，并高效地证明 E 是它们的拓扑直和。把这一观察应用于自同态 \(B_{\lambda} = 1_{E} - \lambda A\) （\(\lambda \in C\) ），他便毫不困难地推出巴拿赫空间上全连续自同态的谱性质，而且这几乎是最终的形式。

关于这些算子的谱的主要结果中，似乎只有那些需要对函数空间概念、特别是对偶概念有更成熟认识的结果才超出了里斯的范围。例如，1920 年代末，T. 希尔德布兰特 [33] 与 J. 绍德尔 [63] 将（在此时已充分确立的巴拿赫空间框架内）证明全连续自同态的伴随自同态是全连续的。

此外，里斯的全连续性概念仍然依赖于序列的使用；巴拿赫 [4] 在其 1932 年的书中摆脱了这一点，转而考虑把每个有界集映为相对紧集的映射。“紧”线性映射这一术语要到二十世纪下半叶才逐渐确立，大概是随着 E. 希尔关于算子半群的书 [34] 而流行开来。

1930 年代提出的若干问题长期悬而未决。在这方面，我们注意到 1973 年对“逼近问题”的解决 [18]，巴拿赫与 S. 马祖尔曾在 1932 与 1938 年使这一问题闻名（见 III, p.~16 的注记 6 与 III, p.~112 的习题 25）。

与此同时，洛莫诺索夫 [40] 证明了：在局部凸空间中，与一个非零紧自同态交换的自同态存在非平凡的不变闭子空间（III, p.~13 的推论 2）。这一结果是对任意连续自同态是否存在此类子空间的问题（“不变子空间问题”）最重大的进展之一。该问题在可分希尔伯特空间的连续自同态情形至今仍未解决；P. 恩夫洛 [19] 在 1987 年构造了第一个不具有非平凡不变闭子空间的巴拿赫空间自同态的例子。
% semantic-fragments: legacy-input-seam
\relax{}
